\documentclass[12pt,a4paper]{article}

% Pacotes básicos
\usepackage[utf8]{inputenc}
\usepackage[T1]{fontenc}
\usepackage[brazil]{babel}
\usepackage{times} % Fonte Times New Roman
\usepackage{geometry}
\usepackage{setspace}
\usepackage{indentfirst}
\usepackage{graphicx}
\usepackage{booktabs}
\usepackage{multirow}
\usepackage{xcolor}
\usepackage{hyperref}
\usepackage{titlesec}
\usepackage{caption}
\usepackage{fancyhdr}
\usepackage{tocloft}
\usepackage[style=abnt, backend=biber]{biblatex}
\addbibresource{bib/referencias.bib} 
% Adequando a NBR 10520:2023
% No arquivo .bib não escreva o nome do autor com caixa alta
\renewcommand*{\UpperOrSCCite}[1]{#1}

\hypersetup{hidelinks}

% Configuração de Margens 
\geometry{
    top=3cm,
    left=3cm,
    bottom=2cm,
    right=2cm
}

% Configuração de Espaçamento e Recuo
\onehalfspacing 
\setlength{\parindent}{1.25cm} 

% Configuração de Citações Longas 
\renewenvironment{quote}
  {\list{}{\leftmargin=4cm \rightmargin=0cm}%
   \item\relax\singlespacing\small}
  {\endlist}

% Formatação dos títulos das seções
\titleformat{\section}{\normalfont\bfseries\MakeUppercase}{\thesection}{1em}{}
\titleformat{\subsection}{\normalfont\bfseries}{\thesubsection}{1em}{}
\titleformat{\subsubsection}{\normalfont}{\thesubsubsection}{1em}{}

% Renomear títulos automáticos para o padrão ABNT
\renewcommand{\listfigurename}{LISTA DE ILUSTRAÇÕES}
\renewcommand{\listtablename}{LISTA DE TABELAS}

% Configuração sumário
\setcounter{tocdepth}{3} % mostra até a seção terciária (1.1.1)
\renewcommand{\cftdotsep}{1}
\renewcommand{\cftsecdotsep}{\cftdotsep}        
\renewcommand{\cftsubsecdotsep}{\cftdotsep}
\renewcommand{\cftsubsubsecdotsep}{\cftdotsep}

% Sem espaço extra entre os itens: o espaçamento vem só do entrelinhas 1,5
\setlength{\cftbeforesecskip}{0pt}
\setlength{\cftbeforesubsecskip}{0pt}
\setlength{\cftbeforesubsubsecskip}{0pt}

% Primárias (
\renewcommand{\cftsecfont}{\normalfont\bfseries}
\renewcommand{\cftsecpagefont}{\normalfont\bfseries}
\renewcommand{\cftsecleader}{\bfseries\cftdotfill{\cftsecdotsep}}
% Secundárias 
\renewcommand{\cftsubsecfont}{\normalfont\bfseries}
\renewcommand{\cftsubsecpagefont}{\normalfont\bfseries}
\renewcommand{\cftsubsecleader}{\bfseries\cftdotfill{\cftsubsecdotsep}}
% Terciárias 
\renewcommand{\cftsubsubsecfont}{\normalfont}
\renewcommand{\cftsubsubsecpagefont}{\normalfont}
\renewcommand{\cftsubsubsecleader}{\normalfont\cftdotfill{\cftsubsubsecdotsep}}

% Todos os níveis alinhados à esquerda 
\cftsetindents{section}{0em}{1.5em}
\cftsetindents{subsection}{0em}{2.5em}
\cftsetindents{subsubsection}{0em}{3.5em}

% Comando que desenha o sumário completo 
\makeatletter
\newcommand{\sumarioabnt}{%
  \begingroup
    \begin{center}\normalfont\bfseries\normalsize SUMÁRIO\end{center}%
    \vspace{\baselineskip}%
    \begin{spacing}{1.5}\@starttoc{toc}\end{spacing}%
  \endgroup
}
\makeatother

% Configuração da lista de ilustrações

\renewcommand{\cftloftitlefont}{\hfill\bfseries\MakeUppercase}
\renewcommand{\cftafterloftitle}{\hfill}

% Configuração da lista de tabelas
\renewcommand{\cftlottitlefont}{\hfill\bfseries\MakeUppercase}
\renewcommand{\cftafterlottitle}{\hfill}


% Configuração de numeração das paginas
\fancypagestyle{plain}{
    \fancyhf{} % Limpa todos os cabeçalhos e rodapés
    \fancyhead[R]{\thepage} % Coloca o número da página à direita (Right)
    \renewcommand{\headrulewidth}{0pt} % Remove a linha horizontal do cabeçalho
    \setlength{\headheight}{14.5pt} % Ajusta a altura para não sobrepor o texto
}

% Remover numeração de página da capa até o sumário
\pagestyle{empty}

\newcommand{\instituicao}{Instituto Federal de Educação, Ciência e Tecnologia do Piauí}
\newcommand{\campus}{Campus Parnaíba}
\newcommand{\curso}{Curso Superior de Tecnologia em Análise e Desenvolvimento de Sistemas - ADS }
\newcommand{\aluno}{Nome do Aluno}
\newcommand{\titulo}{Título do Trabalho}
\newcommand{\disciplina}{Xxxxxxxxxxxx}
\newcommand{\orientador}{Prof. Dr. Xxxxxxxx Xxxxxxxx}
\newcommand{\coorientador}{Prof.ª Dr.ª Xxxxxx Xxx} % Deixe em branco se não houver
\newcommand{\local}{PARNAÍBA -- PIAUÍ}
\newcommand{\ano}{2026}

\begin{document}

% Elementos pré textuais
\begin{titlepage}
    \centering
    \begin{figure}[h]
        \centering
        \includegraphics[width=0.2\textwidth]{images/ifpi_logo.jpg}
    \end{figure}
    \MakeUppercase{\instituicao}\\
    \MakeUppercase{\campus}\\
    \MakeUppercase{\curso}\\
    
    \vspace{5cm}
    
    \MakeUppercase{\aluno}\\
    
    \vspace{4cm}
    
    \MakeUppercase{\textbf{\titulo}\\}
    
    \vfill
    \textbf{
        \MakeUppercase{\local}\\
        \ano
    }
\end{titlepage} 
\begin{titlepage}
    \centering
    \MakeUppercase{\aluno\\}
    
    \vspace{4cm}
    
    \textbf{\MakeUppercase{\titulo}}\\
    
    \vspace{3cm}
    
    \begin{flushright}
        \begin{minipage}{8cm}
            \singlespacing
            Projeto de pesquisa apresentado como exigência para aprovação na disciplina \disciplina do \curso  do \instituicao, \campus. Orientador(a): \orientador. Coorientador(a)(se houver): \coorientador.
        \end{minipage}
    \end{flushright}
    
    \vfill
    
    \local\\
    \ano
\end{titlepage}
\begin{titlepage}
    \centering
    \MakeUppercase{\titulo}
    
    \vspace{4cm}
    
    \MakeUppercase{\textbf{\titulo}}\\
    
    \vspace{3cm}
    
    \begin{flushright}
        \begin{minipage}{8cm}
            \singlespacing
            Projeto de pesquisa apresentado como exigência para aprovação na disciplina \disciplina do \curso do \instituicao, \campus.
        \end{minipage}
    \end{flushright}
    
    \vspace{1cm}
    
    Aprovada em: \_\_\_/\_\_\_/\_\_\_\_\_\_
    
    \vspace{1.5cm}
    
    \textbf{BANCA EXAMINADORA}
    
    \vspace{1cm}
    
    \rule{10cm}{0.4pt}\\
    \orientador (Orientador) \\
    Instituto Federal do Piauí (IFPI)
    
    \vspace{1cm}
    
    \rule{10cm}{0.4pt}\\
    Prof. Dr. Xxxxxxxx Xxxxxxx \\
    Instituição Xxxxxxxx
    
    \vspace{1cm}
    
    \rule{10cm}{0.4pt}\\
    Prof. Dr. Xxxxxxxx Xxxxxxx \\
    Instituição Xxxxxxxx
    
    \vfill
\end{titlepage}
\newpage
\listoffigures
\thispagestyle{empty} % Mantém a página sem numeração
\addcontentsline{toc}{section}{LISTA DE ILUSTRAÇÕES}
\newpage
\listoftables
\thispagestyle{empty}
\addcontentsline{toc}{section}{LISTA DE TABELAS}
\newpage
\begin{center}
    \textbf{LISTA DE ABREVIATURAS}
\end{center}
\addcontentsline{toc}{section}{LISTA DE ABREVIATURAS}
\thispagestyle{empty} % Garante que a página não tenha número

\noindent
\begin{tabular}{@{}ll@{}}
    \textbf{ABNT} & Associação Brasileira de Normas Técnicas \\
    \textbf{ADS}  & Análise e Desenvolvimento de Sistemas \\
    \textbf{CEP}  & Comitê de Ética em Pesquisa \\
    \textbf{IFPI} & Instituto Federal do Piauí \\
    \textbf{TALE} & Termo de Assentimento Livre e Esclarecido \\
    \textbf{TCLE} & Termo de Consentimento Livre e Esclarecido \\
\end{tabular}
\newpage
\begin{center}
    \textbf{LISTA DE SÍMBOLOS}
\end{center}
\addcontentsline{toc}{section}{LISTA DE SÍMBOLOS}

\noindent \% \hspace{1cm} Porcentagem\\
\noindent \copyright \hspace{1cm} Copyright\\
\noindent \textregistered \hspace{1cm} Marca registrada\\
\noindent \$ \hspace{1cm} Dólar\\
\noindent \S \hspace{1cm} Seção\\
\newpage
\thispagestyle{empty}
\sumarioabnt
\thispagestyle{empty}
\newpage

% Elementos textuais
\pagestyle{plain}
\setcounter{page}{9} % Ajusta a numeração para começar na página 9 conforme o sumário do template

\section{INTRODUÇÃO}
É a apresentação do tema e do problema. Deve delimitar o assunto tratado e contextualizá-lo para esclarecer a existência de uma lacuna a ser estudada (problematização). Conforme \textcite[p. 5]{reis_frota} o estudo poderá necessitar de hipótese(s) que são as possibilidades disponíveis que corroboram com as expectativas da solução do problema de pesquisa. Ela(s) podem ser confirmadas ou refutadas no decorrer do estudo.

O texto deve ser digitado usando a fonte Arial ou Times New Roman, espaçamento entrelinhas 1,5 cm, tamanho 12, inclusive para as informações constantes na capa. Para as citações com mais de três linhas, notas de rodapé, paginação, legendas e fontes das ilustrações e das tabelas deve-se usar a fonte no tamanho 10.

O alinhamento do texto deve ser justificado e as referências no final do trabalho devem ser alinhadas à esquerda, com espaçamento simples (1,0) e um espaço simples entre elas, conforme \textcite{abnt6023}. No que tange a contagem das páginas, devem ser feitas a partir da folha de rosto, mas somente numerada a partir da introdução.

Vale salientar que todos os autores citados no texto devem ter a referência incluída em lista no final no trabalho \cite{abnt14724}.

\section{JUSTIFICATIVA}
O texto deve apresentar uma estrutura argumentativa com o objetivo de evidenciar a relevância do tema. Esse tipo de abordagem é fundamental em trabalhos acadêmicos, sobretudo quando se busca financiamento, pois é necessário destacar a importância da pesquisa. Alguns aspectos fundamentais a serem considerados são: qual é a importância da pesquisa para o pesquisador? Por que o tema foi escolhido? Qual é a relevância do estudo para a comunidade acadêmica/científica e para a sociedade?

\begin{figure}[htbp]
    \centering
    \includegraphics[width=0.5\textwidth]{images/ifpi_logo.jpg}
    \caption{Logo do IFPI Parnaíba}
    \label{fig:logoIFPI}
\end{figure}
\section{OBJETIVOS}
Os objetivos da sua pesquisa devem ser divididos em objetivo geral e específicos.

\subsection{OBJETIVO GERAL}
O objeto geral da pesquisa se refere ao tema ou assunto que está sendo investigado. É a questão principal que orienta toda a investigação e define os limites e escopo do estudo. Em outras palavras, o objeto geral da pesquisa é o foco central da investigação, o ponto de partida para a definição dos objetivos específicos e das hipóteses a serem testadas.

\subsection{OBJETIVOS ESPECÍFICOS}
Os objetivos específicos da pesquisa são metas ou etapas intermediárias que precisam ser cumpridas para alcançar o seu objetivo geral. Eles detalham e descrevem as ações que serão realizadas durante a pesquisa, e servem como guia para o pesquisador na condução do estudo. Os objetivos específicos devem ser claros, precisos e mensuráveis, e devem estar diretamente relacionados ao objetivo geral da pesquisa. É comum que cada objetivo específico tenha uma seção específica no trabalho.
\section{REFERENCIAL TEÓRICO}
Destina-se a apresentar as leituras e os fundamentos teóricos que embasam a proposta da pesquisa. O pesquisador deve se atentar para o fato de que citação, conforme \textcite{abnt10520}, de até três linhas não recebe recuo, às maiores que isso deve ser dado um recuo de 4 cm da margem esquerda do texto.

Usar fonte Arial ou Times New Roman, tamanho 10 ou 11, nestes casos: nas citações com mais de três linhas, nas notas de rodapé, na indicação da fonte (autoria) das figuras, quadros e tabelas, bem como nas legendas e nos conteúdos dos quadros e das tabelas. Para a citação com mais de três linhas deixar uma linha em branco antes e outra após a citação direta e recuo de 4 cm da margem esquerda. (Autor, ano, p. xx).

Vale ressaltar que caso a citação seja oriunda de um trabalho em língua estrangeira, cuja tradução foi feita pelo autor do TCC, essa informação deve vir também dentro dos parênteses por meio da expressão “tradução nossa”, conforme exemplificado a seguir:

\begin{quote}
    “citação traduzida citação traduzida citação traduzida citação traduzida citação traduzida.” (Autor, ano, p. xx, tradução nossa).
\end{quote}

Nos casos em que a fonte consultada da citação não possui paginação, a nota de rodapé deve trazer a seguinte informação: A fonte consultada não é paginada. Nesse caso, se enquadram informações extraídas de sites do governo, por exemplo, conforme exemplificado na NBR 10520/2023, item 7.1, exemplo 6:

\begin{quote}
    Exemplo 6: “O Poder Executivo envidará esforços no sentido de antecipar a entrega do Plano previsto no caput deste artigo em pelo menos 15 dias (Brasil, 1999).
    
    Nota: A fonte consultada neste exemplo 6 não é paginada.”
\end{quote}

\subsection{TÍTULO DA SEÇÃO SECUNDÁRIA}
O autor deve nomear as seções de acordo com a sua pesquisa.

\subsubsection{Título da seção terciária}
Texto texto texto texto texto texto texto texto texto texto texto texto texto texto texto texto texto texto texto texto texto texto texto texto texto texto texto texto texto texto texto texto texto texto texto texto texto texto texto texto texto.

\paragraph{Título da seção quaternária}
Texto texto texto texto texto texto texto texto texto texto texto texto texto texto texto texto texto texto texto texto texto.
\section{PROCEDIMENTOS METODOLÓGICOS OU MATERIAL E MÉTODOS}
Aqui são descritos os passos metodológicos percorridos pela pesquisa.

\subsection{MODALIDADE DE PESQUISA}
Quanto ao propósito: pesquisa exploratória, descritiva ou explicativa. Quanto aos métodos empregados: segundo a natureza dos dados (quantitativa, qualitativa ou qualiquantitativa).

Informar o ambiente em que os dados foram coletados (pesquisa de campo, laboratório). O tipo de pesquisa: bibliográfica, documental, experimental, ensaio clínico, estudo de coorte, caso-controle, estudo de caso, etnográfica, pesquisa-ação, pesquisa participante, fenomenológica e mista. \cite{gil2017}

A escolha do procedimento metodológico tem relação direta com a necessidade de atingir os objetivos da pesquisa.

\subsection{INSTRUMENTOS DE COLETA DE DADOS}
Informar quais são os instrumentos de coleta de dados (entrevista, aplicação de questionário, levantamento documental etc.). Como se dará a análise dos dados?

Aqui são descritos os passos metodológicos percorridos pela pesquisa.

\subsection{PARTICIPANTES DA PESQUISA}
Deve informar quem são os partícipes da pesquisa, deixando claros os motivos para a escolha daquele público.

\subsubsection{Aspectos éticos da pesquisa}
Avaliar e apresentar as questões éticas que envolvem a realização da pesquisa. Deve informar se há e quais são os riscos na realização da pesquisa para o público participante. Ex.: riscos ergonômicos, psicológicos, físicos etc.
\section{RESULTADOS ESPERADOS}
Quais são as contribuições que o pesquisador acredita que a realização da pesquisa irá proporcionar para o campo científico e para a sociedade?

\section{ORÇAMENTO}
Deve informar os recursos utilizados para a realização da pesquisa: os recursos humanos, matérias e ou financeiros.

\begin{table}[h]
    \centering
    \caption{Orçamento da Pesquisa}
    
    \label{tab:orcamento}
    \begin{tabular}{|l|l|l|l|}
        \hline
        Itens / despesas & Quantidade & Valor & Quantidade necessária \\ \hline
        01 & & & \\ \hline
        Total & & & \\ \hline
    \end{tabular}
    \vspace{0.2cm}
    
    {\small Fonte: autor (colocar aqui o ano em que está sendo elaborado).}
\end{table}

Informe a fonte de financiamento ou se são recursos próprios. Alerte para as variações de preços que podem ocorrer do espaço de tempo entre o planejamento da pesquisa e a sua execução.
\section{CRONOGRAMA}
O cronograma é a correlação temporal com as atividades previstas. Deve ser construído conforme o planejamento da pesquisa. As etapas podem ser vistas com o(a) orientador (a).

\begin{table}[htbp]
    \centering
    \caption{Cronograma da pesquisa}

    \label{tab:cronograma}
    \resizebox{\textwidth}{!}{
    \begin{tabular}{|l|c|c|c|c|c|c|c|c|c|c|c|c|}
        \hline
        \multirow{2}{*}{ATIVIDADE} & \multicolumn{5}{c|}{2026} & \multicolumn{7}{c|}{2027} \\ \cline{2-13}
        & A & S & O & N & D & J & F & M & A & M & J & J \\ \hline
        Revisão bibliográfica & x & x & x & x & x & x & x & x & x & x & & \\ \hline
        Redação do pré-projeto & x & x & x & & & & & & & & & \\ \hline
        Apresentação do pré-projeto & & x & & & & & & & & & & \\ \hline
        Submissão do projeto ao CEP/IFPI & & & x & & & & & & & & & \\ \hline
        Qualificação do projeto & & & & x & & & & & & & & \\ \hline
        Pesquisa documental & & & & & x & x & & & & & & \\ \hline
        Aplicação de questionário & & & & & & x & x & & & & & \\ \hline
        Análise e tabulação dos dados & & & & & & & & x & x & x & & \\ \hline
    \end{tabular}
    }
\end{table}



\newpage

% Elementos pós textuais
\newpage
\printbibliography[title={REFERÊNCIAS}]
\newpage
\section*{APÊNDICE A - TCLE}
\addcontentsline{toc}{section}{APÊNDICE A - TCLE}
Termo de Consentimento Livre e Esclarecido submetido ao Comitê de Ética em Pesquisa.

{\color{red} Orientações para elaboração}
\newpage
\newpage
\section*{APÊNDICE B - TALE}
\addcontentsline{toc}{section}{APÊNDICE B - TALE}
Termo de Assentimento Livre e Esclarecido para pesquisas com pessoas menores de idade, caso haja.

{\color{red} Orientações para elaboração}

\newpage

\newpage
\section*{ANEXO A – TÍTULO DO ANEXO}
\addcontentsline{toc}{section}{ANEXO A - TÍTULO DO ANEXO}
Os anexos são documentos utilizados pelo autor durante sua pesquisa, mas que não foram criados pelo próprio autor (certidões, leis, jurisprudências, recortes de jornais, etc.) É um elemento opcional, tem função de e/ou comprobatórios do trabalho. \cite{abnt6024}. A palavra “anexo” deve aparecer em letras maiúsculas (cada anexo é descrito por uma letra do alfabeto) seguido de travessão e título.



\end{document}